\subsection{Substitutions}

Let E be an A-algebra. A topology on E is said to be linear if it is invariant under translation and if there exists a fundamental system of neighbourhoods of 0 consisting of ideals of E (Gen.~Top., III, p.~223). The topology on E is then compatible with its A-algebra structure (when A carries the discrete topology). An A-algebra with a linear topology is called a linearly topologized A-algebra.

PROPOSITION 4. --- Let I be a set and E an associative, commutative, unital, linearly topologized separated complete A-algebra.

\begin{enumerate}
\def\labelenumi{(\roman{enumi})}
\tightlist
\item
  Let \(\varphi\) be a continuous homomorphism of A{[}{[}I{]}{]} into E and \(x_{i} = \varphi (\mathbf{X}_{i})\) . Then:
\end{enumerate}

\begin{enumerate}
\def\labelenumi{(\alph{enumi})}
\item
  for all \(i \in \mathbf{I}\) , \(x_{i}^{n}\) tends to 0 as \(n\) tends to \(+\infty\) ;
\item
  if I is infinite, \(x_{i}\) tends to 0 along the filter of complements of finite subsets of I.
\end{enumerate}

\begin{enumerate}
\def\labelenumi{(\roman{enumi})}
\setcounter{enumi}{1}
\tightlist
\item
  Let \(\mathbf{x} = (x_i)_{i\in I}\) be a family of elements of E satisfying \(a\) and \(b\) of (i). Then there exists one and only one unital continuous homomorphism \(\varphi\) of A{[}{[}I{]}{]} into E such that \(\varphi (\mathbf{X}_i) = x_i\) for all \(i\in I\) .
\end{enumerate}

For all \(i \in I\) , \(X_{i}^{n}\) clearly tends to 0 in A{[}{[}I{]}{]} as n tends to \(+\infty\) ; on the other hand when I is infinite, \(X_{i}\) tends to 0 along the filter of complements of finite subsets of I. This proves (i).

Let \((x_{i})_{i\in I}\) be a family of elements of E satisfying conditions \(a)\) and \(b)\) of (i), let \(\psi\) be the homomorphism \(u\mapsto u((x_i)_{i\in I})\) of \(\mathbf{A}[(\mathbf{X}_i)_{i\in I}]\) into E, and let V be a neighbourhood of 0 of E which is an ideal of E. By \(b)\) there exists a finite subset J of I such that \(x_{i}\in \mathbf{V}\) for all \(i\in \mathrm{I - J}\) . Next there exists by \(a)\) an integer \(n\geqslant 0\) such that \(x_{i}^{n}\in \mathbf{V}\) for all \(i\in \mathrm{J}\) . Let \(\beta\) be the element of \(\mathbf{N}^{(I)}\) such that \(\beta_{i} = n - 1\) for \(i\in \mathrm{J}\) and \(\beta_{i} = 0\) for \(i\in \mathrm{I - J}\) . If we define the ideal \(\mathfrak{a}_{\beta}\) of A{[}{[}I{]}{]} as at the beginning of No.~2 (IV, p.~26), then

\[
u \in \mathrm{A} \left[ \left(\mathrm{X} _ {i}\right) _ {i \in I} \right] \cap \mathfrak {a} _ {\beta} \Rightarrow \psi (u) \in \mathrm{V}.
\]

This shows that \(\psi\) is continuous if we equip \(\mathbf{A}[(\mathbf{X}_i)_{i\in I}]\) with the topology induced by that of \(\mathbf{A}[[I]]\) . Since \(\mathbf{E}\) is separated and complete, \(\psi\) extends to a continuous unital homomorphism \(\varphi\) of \(\mathbf{A}[[I]]\) into \(\mathbf{E}\) . We have \(\varphi (\mathbf{X}_i) = \psi (\mathbf{X}_i) = x_i\) for all \(i\in I\) . Finally let \(\varphi^{\prime}\) be a continuous unital homomorphism of \(\mathbf{A}[[I]]\) into \(\mathbf{E}\) such that \(\varphi^{\prime}(\mathbf{X}_{i}) = x_{i}\) . We have \(\varphi^{\prime}(u) = \varphi (u)\) for all \(u\in \mathbf{A}[(\mathbf{X}_i)_{i\in I}]\) , hence \(\varphi^{\prime} = \varphi\) because \(\mathbf{A}[(\mathbf{X}_i)_{i\in I}]\) is dense in \(\mathbf{A}[[I]]\) .

Let us keep the previous notation. If \(u \in \mathbf{A}[[\mathbf{I}]]\) , the image of \(u\) by \(\varphi\) is denoted by \(u(\mathbf{x})\) or \(u((x_i)_{i \in \mathbf{I}})\) (or also \(u(x_1, \ldots, x_n)\) if \(\mathbf{I} = \{1, 2, \ldots, n\}\) ) and is called the element of \(\mathbf{E}\) obtained by substitution of \(x_i\) for \(X_i\) in \(u\) , or the value of \(u\) for the values \(x_i\) of the \(X_i\) or also the value of \(u\) for \(X_i = x_i\) . In particular we have \(u = u((X_i)_{i \in \mathbf{I}})\) .

Let \(E'\) be an associative commutative and unital linearly topologized separated and complete A-algebra. Let \(\lambda\) be a continuous unital homomorphism of E into \(E'\) , and \((x_i)_{i \in I}\) a family of elements of E satisfying conditions \(a)\) and \(b)\) of Prop. 4 (IV, p.~28). The family \((\lambda(x_i))_{i \in I}\) satisfies the same conditions \(a)\) and \(b)\) . For every \(u \in A[[I]]\) we have

\[
\lambda \left(u \left(\left(x _ {i}\right) _ {i \in I}\right)\right) = u \left(\left(\lambda \left(x _ {i}\right)\right) _ {i \in I}\right),\tag{3}
\]

for the mapping \(u \mapsto \lambda(u((x_i)_{i \in I}))\) is a continuous unital homomorphism of \(A[[I]]\) into \(E'\) which transforms \(X_i\) into \(\lambda(x_i)\) for all \(i \in I\) .

If \(J\) and \(K\) are two sets, we denote by \(A_{J,K}\) the set of all families \((g_j)_{j \in J}\) satisfying the following conditions:

\begin{enumerate}
\def\labelenumi{(\roman{enumi})}
\item
  for all \(j \in \mathbf{J}\) , \(g_j\) is an element of \(\mathbf{A}[[\mathbf{K}]]\) whose constant term is nilpotent;
\item
  if \(J\) is infinite, \(g_{j}\) tends to 0 along the filter of complements of finite subsets of \(J\) .
\end{enumerate}

We note that if J is finite, every family of formal power series \((g_{j})_{j\in J}\) without constant term in A{[}{[}K{]}{]} belongs to \(A_{J,K}\) .

Let \((g_j)_{j\in J}\) be in \(\mathbf{A}_{\mathrm{J},\mathbf{K}}\) . By the Cor. of Prop. 3 (IV, p.~28) we have \(\lim_{n\to \infty}g_j^n = 0\) for all \(j\in J\) . Let \(f\in \mathbf{A}[[\mathrm{J}]]\) ; we can substitute \(g_{j}\) for the variable of

index \(j\) in \(f\) and obtain a formal power series \(f((g_j)_{j \in J})\) belonging to A{[}{[}K{]}{]}. Moreover, the mapping \(f \mapsto f((g_j)_{j \in J})\) is a continuous homomorphism of A-algebras A{[}{[}J{]}{]} into A{[}{[}K{]}{]}.

In particular if \(J = \{1, \ldots, p\}\) and \(f \in A[[X_1, \ldots, X_p]]\) , we can substitute for each \(X_j\) a formal power series \(g_j \in A[[K]]\) without constant term; the result of this substitution is written \(f(g_1, \ldots, g_p)\) .

Let \(\mathbf{x} = (x_k)_{k\in K}\) be a family of elements of E satisfying conditions \(a\) and \(b\) ) of Prop. 4 (IV, p.~28). Let us apply (3), taking for \(\lambda\) the homomorphism \(u\mapsto u(\mathbf{x})\) of A{[}{[}K{]}{]} into E; we obtain

\[
f \left(\left(g _ {j}\right) _ {j \in J}\right) (\mathbf {x}) = f \left(\left(g _ {j} (\mathbf {x})\right) _ {j \in J}\right).\tag{4}
\]

Let \(\mathbf{f} = (f_i)_{i\in I}\in (\mathbf{A}[[J]])^I\) and \(g = (g_j)_{j\in J}\in \mathbf{A}_{J,K}\) . We denote by \(\mathbf{f}(\mathbf{g})\) or \(\mathbf{f}\circ \mathbf{g}\) the element \((f_{i}((g_{j})_{j\in J}))_{i\in I}\) of \((\mathbf{A}[[K]])^I\) . If \(\mathbf{f}\in \mathbf{A}_{I,J}\) , we have \(\mathbf{f}\circ \mathbf{g}\in \mathbf{A}_{I,K}\) because the mapping \(f\mapsto f((g_j)_{j\in J})\) of \(\mathbf{A}[[I]]\) into \(\mathbf{A}[[K]]\) is continuous.

Let \(\mathbf{f} \in (\mathbf{A}[[J]])^{\mathrm{I}}\) , \(\mathbf{g} \in \mathbf{A}_{\mathrm{J},\mathrm{K}}\) , \(\mathbf{h} \in \mathbf{A}_{\mathrm{K},\mathrm{L}}\) . Then \(\mathbf{g} \circ \mathbf{h} \in \mathbf{A}_{\mathrm{J},\mathrm{L}}\) and by (4), we have

\[
(\mathbf {f} \circ \mathbf {g}) \circ \mathbf {h} = \mathbf {f} \circ (\mathbf {g} \circ \mathbf {h}).\tag{5}
\]

